\begin{scanpage}{0218}{200}
\[\begin{gathered}
\left.{}+2\bar\psi^2\frac{\ln\dfrac{i}{\gamma ka}}{1+\psi e^{ikL}}e^{2ikL}\right],\\
\bar\psi=\left.\psi_\pm^0\right|_{\vartheta_0=\frac\pi2},
\end{gathered}\tag{32.08}\]
which characterizes the reflected field magnitude with normal irradiation.

Let us also write the expressions for the function $F(\vartheta,\vartheta_0)$ which corresponds to the radar case when the observation and irradiation directions coincide $(\vartheta=\vartheta_0)$
\[\begin{aligned}
F(\vartheta,\vartheta)={}&\frac{i}{\left(\sin\vartheta\ln\dfrac{2i}{\gamma ka\sin\vartheta}\right)^2\cos\vartheta}\Biggl\{\cos^4\frac\vartheta2\ln\frac{i}{\gamma ka\cos^2\dfrac\vartheta2}e^{-2ikz_1\cos\vartheta}\\
&-\sin^4\frac\vartheta2\ln\frac{i}{\gamma ka\sin^2\dfrac\vartheta2}e^{-2ikz_2\cos\vartheta}\\
&+2\left[\sin^4\frac\vartheta2\ln\frac{i}{\gamma ka\sin\dfrac\vartheta2}\psi_-(L)\right.\\
&\qquad\left.-\cos^4\frac\vartheta2\ln\frac{i}{\gamma ka\cos\dfrac\vartheta2}\psi_+(L)\right]e^{-ik(z_1+z_2)\cos\vartheta+ikL}\\
&-2\frac{\cos\vartheta}D\ln\frac{i}{\gamma ka}\psi(L)\psi_+(L)\psi_-(L)e^{-ik(z_1+z_2)\cos\vartheta+3ikL}\\
&+\frac{\cos\vartheta}D\ln\frac{i}{\gamma ka}\bigl[\psi_-^2(L)e^{-2ikz_1\cos\vartheta}+\psi_+^2(L)e^{-2ikz_2\cos\vartheta}\bigr]e^{2ikL}\Biggr\}.
\end{aligned}\tag{32.09}\]
We may show that when $\vartheta=\frac\pi2$ Equation (32.09) leads to Expression (32.08).

The scattering characteristic (32.05) which was found above was obtained by summing all the waves formed with multiple diffraction. Such a method is very graphic, but somewhat lengthy. One may arrive at the same result more quickly if it is assumed that the edge wave diffraction process at the passive vibrator's ends takes place, starting with tertiary action, the same as at the end of a transmitting
\end{scanpage}
